% =========================================================
% T8 — SIMULATION GEOMETRY (Hyper‑Condensed)
% =========================================================

\section{T8 Simulation Geometry}

\paragraph{Overview.}
Simulation Geometry defines the lawful structure of invariant‑restricted
simulation. It provides IRSM, simulation integrity conditions, anchoring
geometry, topological anchors, four‑plane legality validation, projection
planes, mapping planes, invariant validation planes, and CLCP propagation rules
required for stable simulation across DSLO v0.7.

% ---------------------------------------------------------
% T8.1 IRSM
% ---------------------------------------------------------
\subsection{T8.1 IRSM}
Invariant‑Restricted Simulation Mode (IRSM) defines simulation constrained by
curvature invariants, identity boundaries, legality masks, and drift envelopes.
IRSM ensures that simulation behavior remains lawful under load, deformation,
and cross‑layer traversal.

% ---------------------------------------------------------
% T8.2 Simulation Integrity Conditions
% ---------------------------------------------------------
\subsection{T8.2 Simulation Integrity Conditions}
Simulation integrity conditions define the legality and stability requirements
for simulation. Integrity conditions validate curvature continuity, identity
preservation, legality compliance, operator compatibility, and drift coherence
across all simulation transitions.

% ---------------------------------------------------------
% T8.3 Static Anchoring
% ---------------------------------------------------------
\subsection{T8.3 Static Anchoring}
Static anchoring defines topological anchoring for simulation geometry. Anchors
stabilize curvature, identity, and legality under simulation load. Anchoring
prevents drift propagation and collapse during simulation traversal.

% ---------------------------------------------------------
% T8.4 Topological Anchors
% ---------------------------------------------------------
\subsection{T8.4 Topological Anchors}
Topological anchors define invariant anchors across simulation surfaces.
Anchors preserve curvature invariants, identity boundaries, legality masks, and
operator compatibility across simulation layers.

% ---------------------------------------------------------
% T8.5 Four-Plane Legality Check
% ---------------------------------------------------------
\subsection{T8.5 Four-Plane Legality Check}
The Four‑Plane Legality Check validates simulation transitions across four
planes: projection plane, mapping plane, invariant validation plane, and CLCP
plane. The legality check ensures that simulation remains lawful under drift,
load, and deformation.

% ---------------------------------------------------------
% T8.6 Text Projection Plane
% ---------------------------------------------------------
\subsection{T8.6 Text Projection Plane}
The Text Projection Plane defines conceptual vector extraction. Projection
operators map text geometry into simulation geometry while preserving curvature
continuity, identity boundaries, and legality constraints.

% ---------------------------------------------------------
% T8.7 Geometric Mapping Plane
% ---------------------------------------------------------
\subsection{T8.7 Geometric Mapping Plane}
The Geometric Mapping Plane defines manifold mapping. Mapping operators project
simulation geometry across manifold layers while preserving curvature,
identity, legality, and drift coherence.

% ---------------------------------------------------------
% T8.8 Invariant Validation Plane
% ---------------------------------------------------------
\subsection{T8.8 Invariant Validation Plane}
The Invariant Validation Plane validates invariants required for lawful
simulation. Validation ensures that curvature invariants, identity invariants,
legality invariants, and drift invariants remain preserved across simulation
transitions.

% ---------------------------------------------------------
% T8.9 CLCP Plane
% ---------------------------------------------------------
\subsection{T8.9 CLCP Plane}
The CLCP Plane defines cross‑layer constraint propagation. CLCP ensures that
constraints propagate lawfully across manifold layers, preserving curvature,
identity, legality, and drift coherence during simulation traversal.
